% Candidate discussion/conditional proposition; not a standalone document.
% Requires amsmath, amssymb and the submission's theorem/proof environments.
% Labels use tadl:. Full derivations and source audit: deep_learning_model.tex.
% This fragment does not modify jota_submission. The structural hypothesis
% below has not been verified by the currently saved neural norm logs.

\section{A theoretical description of the neural experiments}
\label{tadl:section}

The neural experiments optimize a coupled objective in the hidden matrices
$W_1,\ldots,W_L$ and auxiliary parameters $z$. Applying an LMO separately
to each matrix does not make the objective separable. The appropriate
comparison retains the alignment of the actual joint update with the
target gradient, its directional curvature, and its effective step size.
The CIFAR accuracy results and the language-model results require distinct
geometric descriptions.

\paragraph{Normalized CIFAR training.}
For the weight-normalized airbench variant, a convolution with $m_l$ output
channels is normalized to Frobenius radius $R_l=\sqrt{m_l}$. In the current
implementation normalization precedes the additive update. Write
$N_l(W)=R_lW/\|W\|_{\mathrm F}$ and $Y_l=N_l(W_l)$. If
$D_l=D_l^{\mathrm{tan}}+\alpha_lY_l$, with
$\alpha_l=\langle Y_l,D_l\rangle/R_l^2$, then
\begin{equation}\label{tadl:normalization}
 N_l(Y_l-\gamma_lD_l)
 =N_l\left(Y_l-\frac{\gamma_l}{1-\gamma_l\alpha_l}
                      D_l^{\mathrm{tan}}\right),
 \qquad 1-\gamma_l\alpha_l>0.
\end{equation}
Thus normalized states follow a tangent step and sphere retraction, with a
radial-dependent effective learning rate. Under the idealization of zero
BatchNorm epsilon, positive convolution-weight scalings leave the
training predictor unchanged, and this is a natural geometric description.
Nonzero epsilon and frozen evaluation statistics require separate treatment.
The gradient-norm experiments without weight normalization concern a
different algorithm and do not verify this tangent model.

An ideal rank-$k$ Fanion has Frobenius size $\sqrt{k}$, giving nominal
relative step $\gamma_l\sqrt{k/m_l}$ before tangent projection. Hence
inheriting Muon's learning rate gives Neon an 8--16 times smaller nominal
relative step in these layers. Rank-dependent tuning and actual tangent
displacements are needed to distinguish this effect from direction quality.
For the eight-epoch table, the original per-run normalization flags should
also be checked: a commented Neon call in the current script differs from
the manuscript's declared setting.

\paragraph{Language-model training.}
Use a mean token cross-entropy objective $F_{c_t}(\theta)$, with attention
context rule $c_t$. Hidden matrices receive shape-scaled Fanion updates;
embeddings, head, and scalar parameters receive Adam updates. Decoupled
decay and both hybrid component scalings belong in the actual joint
displacement. RMS normalization of activations does not impose a common
sphere constraint on the hidden weights. Packed-token dependence and the
changing attention window should be represented by noise/bias and objective
drift, respectively. A QKV tensor may contain several independently filtered
matrices while sharing a single Frobenius normalization for its F component.

\paragraph{Directional curvature and useful modes.}
At a fixed state, consider a positive local quadratic model along the
singular directions $P_i=u_iv_i^\top$ of the actual momentum input; for
normalized CIFAR use tangent projections and retraction curvature. Let
$g_i=\langle\operatorname{grad}F,P_i\rangle$ and let $K$ be the modal
curvature Gram matrix. The optimized scalar-step decrease in direction
$\sum_i d_iP_i$ is
\begin{equation}\label{tadl:quality}
 \frac12\mathcal Q(d),\qquad
 \mathcal Q(d)=\frac{(g^\top d)^2}{d^\top Kd},\qquad g^\top d>0.
\end{equation}
The true gradient can be estimated on independent training probes, avoiding
the positive self-alignment of the same stochastic input used by the LMO.
A Gauss--Newton form $J^\top C J$ models curvature in prediction space;
using it as a bound additionally requires control of nonlinear residuals
and cross-layer effects. A global Frobenius smoothness replacement may
discard the structure needed to compare ranks.

\begin{proposition}[Comparable corrections across modes]
\label{tadl:newton}
Suppose the modal quadratic has $K\succ0$ and $g=eK\boldsymbol1$ for
some $e>0$. Then the full ideal Muon direction is a Newton direction in
this model, and every positive-alignment coefficient vector satisfies
\begin{equation}\label{tadl:angle}
 \frac{\mathcal Q(d)}{\mathcal Q(\boldsymbol1)}
 =\frac{(\boldsymbol1^\top Kd)^2}
 {(\boldsymbol1^\top K\boldsymbol1)(d^\top Kd)}\leq1.
\end{equation}
If $K=\operatorname{diag}(\lambda_i)$, this hypothesis means that the
modal errors $g_i/\lambda_i$ are equal. For plain Fanion-$k$ with selected
set $I_k$ it gives
\begin{equation}\label{tadl:mass}
 \frac{\mathcal Q_r}{\mathcal Q_k}
 =\frac{\sum_i\lambda_i}{\sum_{i\in I_k}\lambda_i}
 \geq\frac{\sum_i\lambda_i}{\sum_{i\leq k}\lambda_i^{\downarrow}}.
\end{equation}
The inequality is an equality for the exact-gradient input $H=G$.
More generally, if those positive modal errors lie in
$[e_{\min},e_{\max}]$, then the right-hand side of
\eqref{tadl:mass}, multiplied by $(e_{\min}/e_{\max})^2$, remains a
lower bound on $\mathcal Q_r/\mathcal Q_k$.
\end{proposition}
\begin{proof}
The quadratic's Newton correction is $K^{-1}g=e\boldsymbol1$.
Equation~\eqref{tadl:angle} follows by substitution and Cauchy--Schwarz
in the $K$ inner product. In the diagonal case, for the selected set
$I_k$, $\mathcal Q_k=(\sum_{i\in I_k}\lambda_ie_i)^2/
\sum_{i\in I_k}\lambda_i$. Bound the full score below by
$e_{\min}^2\sum_i\lambda_i$ and the selected score above by
$e_{\max}^2\sum_{i\leq k}\lambda_i^{\downarrow}$.
\end{proof}

The assumption concerns a broad curvature-adjusted correction, not a flat
raw gradient spectrum. Small gradient singular values can reflect small
curvature. A broad correction can therefore favor Muon even when a
Frobenius-gradient surrogate favors Neon. The F-family weights are
$d_i=a\boldsymbol1_{i\leq k}+b\sigma_i(H)/\|H\|_{\mathrm F}$.
For the S family add the sign direction to the local model, preserving
its off-modal curvature and cross terms. Full-rank hybrids close to Muon
in the curvature metric retain a similar score. The condition need not
hold if weak modes are noise or if only a few corrections remain useful.

Finite Newton--Schulz iterations further replace ideal full-polar weights
by $d_i=(p_q\circ\cdots\circ p_1)(\sigma_i(H)/\|H\|_{\mathrm F})$.
The current CIFAR and NanoGPT scripts use three and five repetitions,
respectively, while the Medium script varies the polynomial by iteration.
Keeping these gains gives a model encompassing hard-rank Fanions,
implemented Muon, and both hybrid families. Approximation error and
stochastic momentum alignment cannot be omitted when turning the local
model into a training guarantee.

\paragraph{Finite-horizon guarantees and the evaluation target.}
For a joint step $x_{t+1}=\mathcal R_{x_t}(-\gamma_tV_t)$, suppose
$F_t(x_{t+1})\leq F_t(x_t)-\gamma_tA_t+\gamma_t^2B_t/2$, with
$A_t=\langle\operatorname{grad}F_t,V_t\rangle$. Then
\begin{equation}\label{tadl:finite}
 \sum_{t<T}\gamma_t A_t
 \leq F_0(x_0)-F_{\min}+\frac12\sum_{t<T}\gamma_t^2B_t
       +\sum_{t<T}\big[F_{t+1}(x_{t+1})-F_t(x_{t+1})\big].
\end{equation}
This telescoping bound allows the actual schedules and a changing context.
It does not require that the neural loss be convex or separable.
For CIFAR's $0.2$ label smoothing, the correct loss baseline is target
entropy $H(q)=0.86689391$, and a misclassification costs at least
$\kappa=0.48773025$ above that baseline. Thus for a common predictor and
distribution, $\Pr(\widehat y\ne y)\leq(F-H(q))/\kappa$.
Validation accuracy additionally needs generalization and evaluation-state
control. The saved three gradient norms do not establish these assumptions
or the curvature condition in Proposition~\ref{tadl:newton}; the proposed
mechanism requires direct alignment and curvature measurements.
